\subsection{Rayleigh Quotients and Minimax}

\begin{definition} \label{def:increasing-hermitian-eigenvalue-ordering} \hlt{Increasing Hermitian Eigenvalue Order}\\
For a Hermitian matrix $A\in\operatorname{Mat}(n,n,\mathbb{C})$, write $\lambda_1(A)\leq\lambda_2(A)\leq\cdots\leq\lambda_n(A)$ for its eigenvalues counted with algebraic multiplicity. If only one Hermitian matrix is under discussion, we write $\lambda_i$ instead of $\lambda_i(A)$.
\end{definition}

\begin{theorem} \label{thm:rayleigh-quotient-spectral-interval} \hlt{Rayleigh Quotient on a Spectral Interval}\\
Let $A\in\operatorname{Mat}(n,n,\mathbb{C})$ be Hermitian with increasing eigenvalues $\lambda_1\leq\cdots\leq\lambda_n$ and corresponding orthonormal eigenvectors $\mathbf{u}_1,\ldots,\mathbf{u}_n$. For a unit vector, write $\operatorname{Ray}_A(\mathbf{x})\defeq\mathbf{x}^*A\mathbf{x}$. If $1\leq p\leq q\leq n$ and $\mathbf{x}\in\operatorname{span}(\mathbf{u}_p,\ldots,\mathbf{u}_q)$ is a unit vector, then $\lambda_p\leq \operatorname{Ray}_A(\mathbf{x})\leq \lambda_q$.
\begin{proof}
Write $\mathbf{x}=\sum_{i=p}^{q}c_i\mathbf{u}_i$. Orthonormality gives $\sum_{i=p}^{q}\abs{c_i}^2=1$, and $\mathbf{x}^*A\mathbf{x} =\sum\limits_{i=p}^{q}\lambda_i\abs{c_i}^2$.\\
This is a convex combination of $\lambda_p,\ldots,\lambda_q$, so it lies between the smallest and largest of them.\\
\hlt{Structural Commentary}: In an eigenbasis, the Rayleigh quotient is an eigenvalue average.\\
\hlt{Key Insight}: Spectral coordinates turn a quadratic form into weighted scalar data.
\end{proof}
\end{theorem}

\begin{corollary} \label{cor:rayleigh-extrema-and-diagonal-entry-bounds} \hlt{Rayleigh Extrema and Diagonal Entry Bounds}\\
Let $A\in\operatorname{Mat}(n,n,\mathbb{C})$ be Hermitian.
\begin{enumerate}[label=(\roman*)]
\setlength{\itemsep}{0pt}
\item If $\mathcal{U}_{p,q}=\operatorname{span}(\mathbf{u}_p,\ldots,\mathbf{u}_q)$, then $\min\limits_{\substack{\mathbf{x}\in\mathcal{U}_{p,q}\\ \norm{\mathbf{x}}_2=1}}\mathbf{x}^*A\mathbf{x}=\lambda_p(A)$, and $\max\limits_{\substack{\mathbf{x}\in\mathcal{U}_{p,q}\\ \norm{\mathbf{x}}_2=1}}\mathbf{x}^*A\mathbf{x}=\lambda_q(A).$
\item For every unit vector $\mathbf{x}$, $\lambda_1(A)\leq \mathbf{x}^*A\mathbf{x}\leq\lambda_n(A)$.
\item If $A=[a_{ij}]$, then $\lambda_1(A)\leq a_{ii}\leq\lambda_n(A)$ for $i=1,\ldots,n$.
\end{enumerate}
\begin{proof}
\begin{enumerate}[label=(\roman*)]
\setlength{\itemsep}{0pt}
\item Theorem \ref{thm:rayleigh-quotient-spectral-interval} gives the inequalities on $\mathcal{U}_{p,q}$, with equality at $\mathbf{u}_p$ and $\mathbf{u}_q$.
\item This is (i) with $p=1$ and $q=n$.
\item Apply (ii) to $\mathbf{x}=\mathbf{e}_i$, since $\mathbf{e}_i^*A\mathbf{e}_i=a_{ii}$.
\end{enumerate}
\hlt{Structural Commentary}: Rayleigh extrema on extreme eigendirections; diagonal entries are coordinate quotients.\\
\hlt{Key Insight}: The whole numerical range of a Hermitian matrix sits inside its spectral interval.
\end{proof}
\end{corollary}

\begin{figure}[H]
  \centering
  \begin{tikzpicture}[
  >=Stealth,
  scale=0.88,
  axis/.style={->, gray!65, line width=0.65pt},
  eig/.style={circle, fill=blue!65!black, inner sep=1.7pt},
  rq/.style={circle, fill=red!75!black, inner sep=2.2pt},
  lab/.style={font=\small, align=center, fill=white, inner sep=1pt},
  tiny/.style={font=\small, align=center}
]
  \draw[axis] (-0.4,0) -- (7.4,0);
  \foreach \x/\labx in {0.8/{\lambda_1},2.2/{\lambda_p},5.0/{\lambda_q},6.4/{\lambda_n}} {
    \draw[gray!55] (\x,0.09) -- (\x,-0.09);
    \node[eig] at (\x,0) {};
    \node[tiny, below=4pt] at (\x,0) {$\labx$};
  }
  \draw[green!55!black, line width=1.25pt] (2.2,0.55) -- (5.0,0.55);
  \node[lab, above=4pt] at (3.6,0.55) {$\mathcal{U}_{p,q}$};
  \node[rq] at (3.7,0) {};
  \node[lab, below=16pt] at (3.7,0) {$\operatorname{Ray}_A(\mathbf{x})$};
\end{tikzpicture}

  \caption{For $\mathbf{x}\in\mathcal{U}_{p,q}$, $\operatorname{Ray}_A(\mathbf{x})$ is a convex combination of $\lambda_p,\ldots,\lambda_q$.}
  \label{fig:rayleigh-quotient-convex-combination}
\end{figure}

\begin{theorem} \label{thm:courant-fischer-rayleigh-minmax} \hlt{Rayleigh Minimax Principle}\\
Let $A=A^*\in\operatorname{Mat}(n,n,\mathbb{C})$ have eigenvalues $\lambda_1\leq\cdots\leq\lambda_n$, and set $R_A(\mathbf{x})=\mathbf{x}^*A\mathbf{x}/(\mathbf{x}^*\mathbf{x})$ for $\mathbf{x}\neq\mathbf{0}$.\\
For $1\leq k\leq n$, $\lambda_k =\min\limits_{\substack{\mathcal{S}\subseteq\mathbb{C}^n\\ \dim\mathcal{S}=k}}\; \max\limits_{\mathbf{0}\neq\mathbf{x}\in\mathcal{S}} R_A(\mathbf{x}) =\max\limits_{\substack{\mathcal{S}\subseteq\mathbb{C}^n\\ \dim\mathcal{S}=n-k+1}}\; \min\limits_{\mathbf{0}\neq\mathbf{x}\in\mathcal{S}} R_A(\mathbf{x})$.
\begin{proof}
Let $\mathbf{u}_1,\ldots,\mathbf{u}_n$ be an orthonormal eigenbasis in increasing order.\\
Set $E_k=\operatorname{span}(\mathbf{u}_1,\ldots,\mathbf{u}_k)$; $F_k=\operatorname{span}(\mathbf{u}_k,\ldots,\mathbf{u}_n)$.\\
By Theorem \ref{thm:rayleigh-quotient-spectral-interval}, $\max_{\mathbf{0}\neq\mathbf{x}\in E_k}R_A(\mathbf{x})=\lambda_k$ and $\min_{\mathbf{0}\neq\mathbf{x}\in F_k}R_A(\mathbf{x})=\lambda_k$.\\
If $\dim\mathcal{S}=k$, then $\mathcal{S}\cap F_k\neq\{\mathbf{0}\}$, so $\max_{\mathbf{0}\neq\mathbf{x}\in\mathcal{S}}R_A(\mathbf{x})\geq\lambda_k$; equality is attained by $E_k$.\\
If $\dim\mathcal{T}=n-k+1$, then $\mathcal{T}\cap E_k\neq\{\mathbf{0}\}$, so $\min_{\mathbf{0}\neq\mathbf{x}\in\mathcal{T}}R_A(\mathbf{x})\leq\lambda_k$; equality is attained by $F_k$.\\
\hlt{Structural Commentary}: Dimension forces intersections with spectral flags.\\
\hlt{Key Insight}: Minimax is Rayleigh's quotient plus dimension counting.
\end{proof}
\end{theorem}

\begin{method} \label{meth:rayleigh-eigenbasis-method} \hlt{Rayleigh Eigenbasis Method}\\
To bound a Hermitian quadratic form:
\begin{enumerate}[label=(\roman*)]
\setlength{\itemsep}{0pt}
\item diagonalise $A$ with an orthonormal eigenbasis;
\item expand the unit vector in that basis;
\item rewrite $\mathbf{x}^*A\mathbf{x}$ as $\sum_i\lambda_i\abs{c_i}^2$;
\item use the eigenvalue range allowed by the support of $\mathbf{x}$.
\end{enumerate}
\end{method}

